\chapter{Connectivité on-chain}\label{ch:nw:on-chain-connectivity}

Une transaction envoyée à un nœud public d'une blockchain atteint le reste du réseau par gossip, saut après saut, et est visible de tous sur le
chemin~; celle envoyée à un \omterm{def:nw:on-chain-connectivity:engine}{moteur de blocs} situé dans la région du prochain producteur de blocs contourne entièrement le gossip. Sur une chaîne dont les leaders changent toutes les
quelques centaines de millisecondes, et dont le logiciel attribue à chaque leader quatre créneaux consécutifs, la région d'envoi relève d'une question d'ordonnancement~: la
réponse change avec le \omterm{def:nw:on-chain-connectivity:leader}{calendrier des leaders}, et ce calendrier est publié à l'avance.

Le livre~3 a décrit les acteurs du trading on-chain (nœuds, validateurs, constructeurs de blocs, relais, searchers). Ce chapitre porte sur leur réseau~: comment
les nœuds se connectent entre pairs et diffusent par gossip, ce que changent les réseaux de propagation et les relais privés, où se trouvent les \omterm{def:nw:on-chain-connectivity:engine}{moteurs de blocs}, comment un \omterm{def:nw:on-chain-connectivity:leader}{calendrier des leaders} se traduit en
choix d'endpoint, et ce que vendent les \omterm{def:nw:on-chain-connectivity:engine}{fournisseurs de nœuds hébergés}. La documentation est citée et datée~; la propagation fait l'objet d'une simulation étiquetée sur
la carte du livre.

\section{Exploiter et interconnecter des nœuds}

\begin{definition}[Protocole de bavardage (gossip)]\label{def:nw:on-chain-connectivity:gossip}
Un \emph{protocole de bavardage}\index{protocole de bavardage} (gossip protocol) diffuse un message sur un réseau pair à pair en faisant en sorte que chaque nœud qui le reçoit pour la
première fois le transmette (ou l'annonce) à certains ou à tous ses pairs, de sorte qu'il atteigne chaque nœud en un nombre de sauts croissant avec le logarithme de
la taille du réseau.
\end{definition}

Un nœud rejoint le réseau d'une chaîne en se connectant à des pairs trouvés par le protocole de découverte de la chaîne~; dans le modèle de ce chapitre, ils sont choisis au
hasard, sans égard à la géographie. Chaque saut coûte le délai du lien, auquel s'ajoute la validation du message par le nœud avant qu'il ne le transmette. La conséquence a été mesurée très tôt~:
Decker et Wattenhofer ont établi en 2013 qu'un bloc Bitcoin atteignait le nœud médian au bout de 6.5 secondes et qu'après 40 secondes 5~\% des nœuds n'avaient
pas encore vu ce bloc. Des messages plus petits et une validation plus rapide raccourcissent chaque saut, mais pas la structure~: une transaction atteint
l'autre bout du monde par une suite de sauts entre des pairs qu'elle n'a pas choisis.

\begin{proposition}[Temps de diffusion par gossip]\label{prop:nw:on-chain-connectivity:gossip}
Sur un réseau dont chaque nœud transmet à tous ses pairs dès la première réception, un message atteint chaque nœud au bout de la longueur de son plus court chemin
depuis la source, la longueur d'une liaison étant son délai augmenté du traitement par le nœud récepteur~; transmettre à un sous-ensemble et annoncer aux autres, qui
la récupèrent (pull), allonge d'un aller-retour les liaisons annoncées et ne peut qu'allonger le délai d'arrivée.
\end{proposition}

\begin{proof}
L'inondation explore tous les chemins en parallèle, si bien que chaque nœud reçoit le message en premier par son plus court chemin (argument de Dijkstra). Remplacer certaines
transmissions (push) par des annonces ajoute un aller-retour à la longueur de ces liaisons, et les plus courts chemins ne peuvent pas devenir plus courts quand les longueurs croissent.
\end{proof}

\omcode{code/firm/chainnet/firm_chainnet.py}{54}{71}{Gossip comme plus courts chemins~: inondation, et diffusion en racine carrée avec annonces et récupérations.}\label{lst:nw:on-chain-connectivity:spread}

Le modèle répartit 600 nœuds uniformément dans les huit villes listées dans l'\cref{dat:nw:on-chain-connectivity:engines}, relie chacun à huit pairs tirés au hasard, et
facture à chaque liaison son plancher de fibre multiplié par 1.5, plus \qty{5}{\milli\second} de traitement (hypothèses). Une transaction émise depuis un nœud de Tokyo atteint
la moitié des nœuds après \qty{85.0}{\milli\second} et les neuf dixièmes après \qty{99.5}{\milli\second} par inondation~; avec une diffusion en racine carrée, qui
économise la bande passante, après 113.9 et \qty{149.0}{\milli\second}. Dans le modèle, le graphe aléatoire offre des chemins proches des liaisons directes,
si bien que ce qu'ajoute le gossip, ce sont ses sauts supplémentaires~: trois à huit, de 15 à \qty{40}{\milli\second} de traitement, contre celui du moteur,
\qty{1}{\milli\second} (hypothèse). Ce sont les queues qui diffèrent le plus~: le dernier nœud la reçoit après \qty{158.5}{\milli\second} par inondation et
\qty{326.5}{\milli\second} avec une diffusion en racine carrée.

\begin{omfigure}
\begin{tikzpicture}
\begin{axis}[width=11cm,height=5.6cm,xmin=0,xmax=180,ymin=0,ymax=1.02,xlabel={temps écoulé depuis le départ de la transaction de Tokyo (ms)},
  ylabel={part des nœuds atteints},tick label style={font=\footnotesize},label style={font=\small},grid=major,grid style={gray!25},
  legend pos=north west,legend style={font=\footnotesize},table/col sep=comma]
\addplot[omDef,thick,const plot] table[x=ms,y=flood]{figdata/networks/21-on-chain-connectivity/propagation.csv};
\addplot[omThm,thick,dashed,const plot] table[x=ms,y=sqrt]{figdata/networks/21-on-chain-connectivity/propagation.csv};
\addplot[gray,densely dotted] coordinates {(0,0.5) (180,0.5)};
\addplot[gray,densely dotted] coordinates {(0,0.9) (180,0.9)};
\legend{inondation,diffusion en racine carrée}
\end{axis}
\end{tikzpicture}

\omcaption{Simulation~: part des 600 nœuds répartis dans huit villes, dotés chacun de huit pairs aléatoires, ayant reçu une transaction envoyée d'un nœud de Tokyo, en fonction
du temps. L'inondation atteint la moitié à \qty{85.0}{\milli\second} et les neuf dixièmes à \qty{99.5}{\milli\second}~; la diffusion en racine carrée atteint 113.9 et
\qty{149.0}{\milli\second}. Données~: \texttt{fig\_chain.py}, \texttt{firm\_chainnet.spread()}.}\label{fig:nw:on-chain-connectivity:propagation}
\end{omfigure}

\begin{omfigure}
\begin{tikzpicture}
\begin{axis}[xbar,width=11cm,height=6.2cm,bar width=6pt,xmin=0,xmax=110,ytick=data,
  yticklabels={Amsterdam,Dublin,Francfort,Londres,New York,Salt Lake City,Singapour,Tokyo},y dir=reverse,
  xlabel={arrivée chez un leader de la ville (ms), depuis Tokyo},tick label style={font=\footnotesize},label style={font=\small},
  grid=major,grid style={gray!25},legend style={font=\footnotesize,at={(0.5,-0.3)},anchor=north,legend columns=2},
  table/col sep=comma]
\addplot[fill=omThm!45,draw=omThm] table[y expr=\coordindex,x=gossip]{figdata/networks/21-on-chain-connectivity/leaders.csv};
\addplot[fill=omDef!55,draw=omDef] table[y expr=\coordindex,x=direct]{figdata/networks/21-on-chain-connectivity/leaders.csv};
\legend{par gossip,par le moteur de blocs le plus proche}
\end{axis}
\end{tikzpicture}

\omcaption{Simulation~: moment où une transaction envoyée depuis Tokyo atteint le leader suivant, par ville du leader~: par le graphe de gossip depuis un nœud public
à Tokyo (médiane sur les nœuds de la ville), et directement par le \omterm{def:nw:on-chain-connectivity:engine}{moteur de blocs} publié le plus proche du leader. Le gossip ajoute 14 à
\qty{39}{\milli\second}. Données~: \texttt{fig\_chain.py}, \texttt{nw\_chain.to\_leader()}.}\label{fig:nw:on-chain-connectivity:leaders}
\end{omfigure}

\section{Réseaux de propagation des transactions et relais privés}

\begin{definition}[Réseau de propagation des transactions]\label{def:nw:on-chain-connectivity:tpn}
Un \emph{réseau de propagation des transactions}\index{réseau de propagation des transactions} (transaction propagation network) est un réseau de relais exploité en privé, interconnecté avec de nombreux nœuds
et validateurs d'une chaîne, qui fait circuler transactions et blocs entre eux plus vite et avec moins de variance que le gossip propre à la chaîne sur
l'internet public.
\end{definition}

Deux services modifient le trajet d'une transaction. Un réseau de propagation, tel que le Blockchain Distribution Network de bloXroute, remplace les sauts publics aléatoires
par ses propres relais, que sa documentation décrit comme réduisant «~la latence et la variance~» dans la propagation. Un relais privé modifie qui voit
la transaction~: Flashbots Protect, sur Ethereum, l'envoie vers un mempool privé où elle est «~cachée aux bots de frontrunning et de sandwich~» jusqu'à ce qu'un
constructeur de blocs l'inclue. La première est un produit de latence~; la seconde un produit de visibilité~; un searcher peut avoir besoin des deux (le flux d'ordres privé du livre~3).

\section{Endpoints des constructeurs de blocs et des moteurs de blocs par région}

\begin{definition}[Moteur de blocs, fournisseur de nœuds hébergés]\label{def:nw:on-chain-connectivity:engine}
Un \emph{moteur de blocs}\index{moteur de blocs} (block engine) est un service exploité à proximité des validateurs d'une chaîne, qui accepte les transactions et les bundles des traders et
les transmet directement au producteur de blocs en cours ou au suivant, souvent à travers une enchère d'inclusion. Un
\emph{fournisseur de nœuds hébergés}\index{fournisseur de nœuds hébergés} (hosted node provider) fait tourner les nœuds d'une chaîne et vend l'accès à leurs endpoints RPC, avec des limites de requêtes par forfait, si bien qu'une société n'a pas à les exploiter elle-même.
\end{definition}

\begin{dated}{2026-09}{Moteurs de blocs et relais, d'après leur documentation}\label{dat:nw:on-chain-connectivity:engines}
\begin{itemize}
\item Jito (Solana)~: \omterm{def:nw:on-chain-connectivity:engine}{moteurs de blocs} du mainnet à Amsterdam, Dublin, Francfort, Londres, New York, Salt Lake City, Singapour et Tokyo, chacun avec son propre
  endpoint~; sa \texttt{sendTransaction} transmet une transaction directement au validateur~; les bundles exigent un pourboire d'au moins 1\,000 lamports.
\item Flashbots Protect (Ethereum)~: un mempool privé à l'abri des bots de frontrunning et de sandwich, avec des remboursements.
\item bloXroute~: un Blockchain Distribution Network de relais pour Solana, BNB Chain, Base, Ethereum, Hyperliquid et d'autres chaînes.
\end{itemize}
\end{dated}

\begin{omfigure}
\begin{tikzpicture}[font=\small,node distance=6mm,
  nd/.style={circle,draw=omThm,fill=omThm!15,minimum size=5mm,inner sep=0pt},
  box/.style={draw=omDef,fill=omDef!12,rounded corners,align=center,minimum height=8mm,inner sep=3pt}]
\node[box] (firm) at (0,0) {société\\(Tokyo)};
\node[nd] (a) at (2.2,1.3) {}; \node[nd] (b) at (3.8,1.9) {}; \node[nd] (c) at (5.4,1.2) {};
\node[nd] (d) at (7.0,1.8) {}; \node[nd] (e) at (4.6,0.3) {};
\node[box] (lead) at (9.4,0) {prochain leader\\(Francfort)};
\node[box] (eng) at (7.4,-1.3) {moteur de blocs\\(Francfort)};
\draw[omThm,->] (firm) -- (a); \draw[omThm,->] (a) -- (b); \draw[omThm,->] (b) -- (c); \draw[omThm,->] (c) -- (d);
\draw[omThm,->] (d) -- (lead); \draw[omThm!50,->] (a) -- (e); \draw[omThm!50,->] (e) -- (c);
\draw[omDef,->,thick] (firm) |- (eng); \draw[omDef,->,thick] (eng) -- (lead);
\node[text=omThm,font=\footnotesize] at (4.6,2.5) {gossip public : chaque saut le voit et y ajoute son traitement};
\node[text=omDef,font=\footnotesize] at (2.6,-2.0) {route propre jusqu'au moteur ; un dernier saut court};
\end{tikzpicture}

\omcaption{Deux chemins vers le prochain leader~: confier la transaction à un nœud public et laisser le gossip la transporter de proche en proche, visible de chaque nœud~; ou l'envoyer
par la route propre de la société jusqu'au \omterm{def:nw:on-chain-connectivity:engine}{moteur de blocs} le plus proche du leader, qui la transmet directement. Schéma.}\label{fig:nw:on-chain-connectivity:paths}
\end{omfigure}

Un \omterm{def:nw:on-chain-connectivity:engine}{moteur de blocs} (block engine) est un \omterm{def:nw:the-european-map:pop}{point de présence} (chapitre~11) pour les producteurs de la chaîne. Envoyer au moteur le plus proche du producteur place la partie longue distance du trajet
sur la route propre de la société et le dernier saut sur celle, courte, du moteur, au lieu de passer par une douzaine de sauts publics aléatoires.

\section{Proximité des validateurs et calendrier des leaders}

\begin{definition}[Calendrier des leaders]\label{def:nw:on-chain-connectivity:leader}
Un \emph{calendrier des leaders}\index{calendrier des leaders} est l'affectation publiée par une chaîne des créneaux de production de blocs à venir aux validateurs, calculée à
l'avance pour une époque, ce qui permet à chacun de savoir quel validateur produira chaque bloc à venir et, si sa position est connue, où envoyer
les transactions.
\end{definition}

\begin{dated}{2026-09}{Créneaux et leaders, d'après les spécifications et le code des chaînes}\label{dat:nw:on-chain-connectivity:chains}
La configuration du réseau principal d'Ethereum fixe des créneaux de 12~secondes. Solana calcule son \omterm{def:nw:on-chain-connectivity:leader}{calendrier des leaders} une époque à l'avance~; la durée cible par défaut d'un créneau dans le SDK de Solana
est de \qty{300}{\milli\second} (il définit aussi 400 et moins) et la cible effective est modifiée dynamiquement, si bien que les clients ne doivent pas supposer
la valeur par défaut~; le validateur Agave attribue à chaque leader quatre créneaux consécutifs.
\end{dated}

Avec des créneaux de 12~secondes, la géographie ne compte pour Ethereum qu'aux extrémités d'un créneau. Avec des créneaux de quelques centaines de millisecondes et quatre d'affilée par
leader, elle compte pour chaque transaction sur Solana~: la position du prochain leader est connue à l'avance, et une transaction qui l'atteint après ses
quatre créneaux attend le leader suivant, peut-être sur un autre continent.

\omcode{code/networks/21-on-chain-connectivity/python/nw_chain.py}{36}{44}{Arrivée au prochain leader~: par le graphe de gossip, et directement par
le moteur publié le plus proche du leader.}\label{lst:nw:on-chain-connectivity:leader}

\begin{omfigure}
\begin{tikzpicture}
\begin{axis}[width=11cm,height=5.6cm,xmin=0,xmax=150,ymin=0,ymax=1.05,xlabel={temps restant dans la fenêtre du leader au moment de l'envoi de la transaction (ms)},
  ylabel={part des leaders atteints},tick label style={font=\footnotesize},label style={font=\small},grid=major,grid style={gray!25},
  legend pos=north west,legend style={font=\footnotesize},table/col sep=comma]
\addplot[omThm,thick,const plot] table[x=deadline,y=gossip]{figdata/networks/21-on-chain-connectivity/timely.csv};
\addplot[omDef,thick,const plot] table[x=deadline,y=direct]{figdata/networks/21-on-chain-connectivity/timely.csv};
\legend{par le gossip,moteur conscient du leader}
\end{axis}
\end{tikzpicture}

\omcaption{Simulation~: la part des prochains leaders, placés tour à tour dans chacune des huit villes des moteurs, atteints à temps depuis Tokyo, en fonction du temps restant
dans la fenêtre du leader. Envoyer au moteur le plus proche du leader décale la courbe vers la gauche de 14 à \qty{39}{\milli\second}. Données~: \texttt{fig\_chain.py},
\texttt{nw\_chain.timely\_rates()}.}\label{fig:nw:on-chain-connectivity:timely}
\end{omfigure}

Moyenne prise sur les huit villes des leaders et sur un temps restant uniforme entre zéro et \qty{150}{\milli\second}, le sélecteur conscient du leader fait parvenir à temps 61~\% des
transactions contre 48~\% par le gossip~; sur \qty{300}{\milli\second}, 80.5~\% contre 74.0~\%. Pour un leader à Francfort, le
problème du chapitre, le chemin direct prend \qty{69.5}{\milli\second} et le gossip \qty{83.5}{\milli\second}.

\section{Fournisseurs de nœuds hébergés}

Faire tourner correctement les nœuds d'une chaîne demande un travail important~: stockage, bande passante, mises à jour et peering. Les fournisseurs hébergés
vendent à la place les endpoints de leurs nœuds, selon une formule et une limite de requêtes. Pour le trading, la question est celle des chapitres~16 à~18~: où se trouvent
les nœuds du fournisseur, s'il s'agit de ceux par lesquels les transactions de la société doivent entrer, et comment leurs limites sont gouvernées (les limites de débit du livre~3). Un
endpoint de fournisseur est un lieu commode pour lire la chaîne et un mauvais lieu pour y rivaliser.

\begin{method}[Faire parvenir une transaction au bloc suivant]\label{met:nw:on-chain-connectivity:send}
\begin{enumerate}
\item Lire le temps de créneau et le \omterm{def:nw:on-chain-connectivity:leader}{calendrier des leaders} de la chaîne~; situer les prochains leaders dans les régions où leur position est connue ou peut être mesurée
  (triangulation du chapitre~17).
\item Garder des connexions ouvertes vers les \omterm{def:nw:on-chain-connectivity:engine}{moteurs de blocs} ou les constructeurs de blocs dans chaque région utilisée par les leaders~; choisir pour chaque transaction celui le plus proche du prochain
  leader avec assez de temps restant.
\item Utiliser un réseau de propagation ou un relais privé lorsque la visibilité ou la variance comptent plus que les dernières millisecondes.
\item Faire tourner ses propres nœuds pour lire l'état critique au trading~; utiliser les endpoints hébergés pour le reste, dans leurs limites.
\item Mesurer les taux d'atterrissage par région et par leader~; le calendrier et les emplacements des moteurs changent.
\end{enumerate}
\end{method}

\section{Tutoriel~: gossip et prochain leader}\label{tut:nw:on-chain-connectivity:tut}

\begin{tutorial}
\textbf{Objectif.} Simuler le gossip sur un graphe de pairs au-dessus de la carte du livre, et le comparer à l'envoi au \omterm{def:nw:on-chain-connectivity:engine}{moteur de blocs} le plus proche du prochain leader.
\textbf{État final~:} \cref{fig:nw:on-chain-connectivity:propagation,fig:nw:on-chain-connectivity:leaders,fig:nw:on-chain-connectivity:timely}.
\begin{tutsteps}
\item \textbf{Moteurs.} \texttt{firm\_chainnet.load\_engines} lit les endpoints régionaux datés.
\item \textbf{Graphe.} \texttt{Graph} place 600 nœuds dans les villes des moteurs, avec huit pairs chacun~; \texttt{spread} (\cref{lst:nw:on-chain-connectivity:spread})
  donne le temps d'arrivée de chaque nœud.
\item \textbf{Leaders.} \texttt{nw\_chain.to\_leader} (\cref{lst:nw:on-chain-connectivity:leader}) compare le gossip avec le moteur le plus proche~;
  \texttt{timely\_rates} et \texttt{mean\_timely} le convertissent en parts d'atterrissage.
\end{tutsteps}
\textbf{Pour aller plus loin.} Rendre le graphe sensible à la géographie (pairs choisis parmi les nœuds les plus proches) et voir ce que le gossip y gagne~; ajouter la file d'attente propre aux moteurs
et leurs pourboires.
\end{tutorial}

\section{Construction~: le modèle de réseau de la chaîne}\label{bld:nw:on-chain-connectivity:chainnet}

\begin{build}
\bfield{Rôle} Comment les transactions circulent sur le réseau d'une chaîne et où les envoyer~: les lignes on-chain du plan du chapitre~29.
\bfield{Interface} \texttt{firm\_chainnet}~: \texttt{Graph}, \texttt{spread}, \texttt{reach\_time}, \texttt{Engine}, \texttt{load\_engines},
\texttt{one\_way\_ms}, \texttt{choose\_engine}, \texttt{timely}.
\bfield{Règles} Les endpoints sont des lignes datées avec leurs sources~; les délais des liaisons sont le plancher de la fibre multiplié par un \omterm{def:nw:the-north-american-map:factor}{facteur de tracé} déclaré, plus un traitement déclaré~; chaque résultat
est étiqueté comme une simulation.
\bfield{Tests d'acceptation} \texttt{code/firm/chainnet/tests/}~: les données des moteurs, l'inondation sur un graphe en ligne vérifiée à la main, la diffusion en racine carrée jamais plus rapide
que l'inondation, et le contrôle du sélecteur et de l'échéance.
\bfield{Pour aller plus loin} Lire un \omterm{def:nw:on-chain-connectivity:leader}{calendrier des leaders} en direct et les positions mesurées des validateurs~; modéliser les enchères des moteurs.
\end{build}

\begin{omsources}
\item C.~Decker et R.~Wattenhofer, «~Information propagation in the Bitcoin network~», IEEE P2P 2013.
\item Ethereum consensus specifications (configuration mainnet)~; Anza documentation (rotation des leaders), Solana SDK et code source d'Agave.
\item Jito Labs, low-latency transaction send~; Flashbots Protect~; bloXroute documentation.
\end{omsources}

\section{Exercices}

\begin{exercise}[$\star$]\label{exo:nw:on-chain-connectivity:1}
Calculer pendant combien de millisecondes un leader Solana garde la chaîne s'il dispose de quatre slots consécutifs à la cible par défaut du SDK.
\end{exercise}

\begin{exercise}[$\star$]\label{exo:nw:on-chain-connectivity:2}
Expliquer la différence entre un \omterm{def:nw:on-chain-connectivity:tpn}{réseau de propagation des transactions} et un relais privé.
\end{exercise}

\begin{exercise}[$\star$]\label{exo:nw:on-chain-connectivity:3}
Expliquer pourquoi le gossip atteint plus tard un nœud situé dans la ville même de l'émetteur qu'un message direct le ferait.
\end{exercise}

\begin{exercise}[$\star\star$]\label{exo:nw:on-chain-connectivity:4}
En utilisant la \cref{prop:nw:on-chain-connectivity:gossip}, expliquer pourquoi la diffusion en racine carrée n'est jamais plus rapide que l'inondation, et pourquoi l'utiliser.
\end{exercise}

\begin{exercise}[$\star\star$]\label{exo:nw:on-chain-connectivity:5}
Le prochain leader se trouve à Singapour, et la machine à Tokyo~: indiquer quel moteur utiliser et quel gain le modèle annonce.
\end{exercise}

\begin{exercise}[$\star\star$]\label{exo:nw:on-chain-connectivity:6}
Expliquer pourquoi le choix de la région compte peu sur une chaîne aux slots de 12~secondes.
\end{exercise}

\begin{exercise}[$\star\star\star$]\label{exo:nw:on-chain-connectivity:7}
\emph{Programmation.} Avec \texttt{firm.chainnet}, calculer de combien plus vite l'inondation atteint les neuf dixièmes des nœuds si chaque nœud a 16 pairs au lieu de 8.
\end{exercise}

\begin{exercise}[$\star\star\star$]\label{exo:nw:on-chain-connectivity:8}
\emph{Trouver la faille.} «~Nous envoyons chaque transaction au \omterm{def:nw:on-chain-connectivity:engine}{moteur de blocs} le plus proche de nos serveurs, donc nous empruntons toujours le chemin le plus rapide.~»
\end{exercise}

\section{Problème~: Le leader est à Francfort}

\begin{problem}[{Problème du week-end --- gossip contre le moteur attentif au leader}]\label{pb:nw:on-chain-connectivity:1}
Une société de trading à Tokyo envoie des transactions sur une chaîne dont le prochain leader se trouve à Francfort. Elle peut les confier à un nœud public à Tokyo et laisser le gossip les acheminer,
ou les envoyer au \omterm{def:nw:on-chain-connectivity:engine}{moteur de blocs} le plus proche du leader. Utiliser le modèle de réseau du chapitre.

\medskip\noindent\textbf{Partie I --- Gossip.}
\begin{enumerate}
\item Combien de temps l'inondation met-elle pour atteindre la moitié et les neuf dixièmes des nœuds à partir de Tokyo~?
\item Et la diffusion en racine carrée~?
\item Quand la transaction atteint-elle un leader à Francfort par le gossip~?
\item Pourquoi cela dure-t-il plus longtemps que le chemin direct~?
\end{enumerate}

\medskip\noindent\textbf{Partie II --- Le moteur.}
\begin{enumerate}[resume]
\item Quel moteur est le plus proche du leader, et quand la transaction atteint-elle le leader par son intermédiaire~?
\item Quel est le plancher de la fibre dans un sens entre Tokyo--Francfort, et quel \omterm{def:nw:the-north-american-map:factor}{facteur de tracé} le modèle retient-il~?
\item Quelle part du chemin direct représente le temps propre du moteur~?
\item Que voit le moteur que le gossip aurait révélé à tous~?
\end{enumerate}

\medskip\noindent\textbf{Partie III --- Taux d'atteinte.}
\begin{enumerate}[resume]
\item Le temps restant étant uniforme entre zéro et \qty{150}{\milli\second}, quelle part des leaders des huit villes chaque méthode atteint-elle~?
\item Et jusqu'à \qty{300}{\milli\second}~?
\item Pourquoi le gain diminue-t-il quand la fenêtre s'élargit~?
\item Comment quatre slots consécutifs par leader modifieraient-ils la planification de la société~?
\end{enumerate}

\medskip\noindent\textbf{Partie IV --- Le verdict.}
\begin{enumerate}[resume]
\item Énoncer le \emph{résultat nommé}~: le temps de propagation jusqu'à la moitié et aux neuf dixièmes du réseau, et le gain d'arrivée à temps obtenu en envoyant au
  \omterm{def:nw:on-chain-connectivity:engine}{moteur de blocs} le plus proche du prochain leader, comparé à une région fixe.
\item Quelles entrées relèvent de la documentation et lesquelles des hypothèses~?
\item Que changerait un graphe de pairs tenant compte de la géographie~?
\item Où les machines de la société doivent-elles se trouver pour cette chaîne~?
\item Qu'apporte un réseau de propagation en plus du moteur~?
\item Comment la société doit-elle vérifier que son modèle correspond à la réalité~?
\item Où cela se range-t-il dans le plan du chapitre~29~?
\item En une phrase~: pourquoi la région est-elle une question de calendrier~?
\end{enumerate}
\end{problem}

\section{Questions d'entretien}

\begin{interviewq}[$\star$ \iqroles{développeur}]\label{iq:nw:on-chain-connectivity:1}
Comment une transaction se propage-t-elle sur le réseau pair à pair d'une blockchain, et pourquoi est-ce lent~?
\end{interviewq}

\begin{interviewq}[$\star\star$ \iqroles{développeur}]\label{iq:nw:on-chain-connectivity:2}
Qu'est-ce qu'un \omterm{def:nw:on-chain-connectivity:engine}{moteur de blocs}, et pourquoi en existe-t-il plusieurs dans le monde~?
\end{interviewq}

\begin{interviewq}[$\star\star$ \iqroles{développeur, chercheur}]\label{iq:nw:on-chain-connectivity:3}
Comment utiliseriez-vous le \omterm{def:nw:on-chain-connectivity:leader}{calendrier des leaders} d'une chaîne pour acheminer les transactions~?
\end{interviewq}

\begin{interviewq}[$\star\star$ \iqroles{développeur}]\label{iq:nw:on-chain-connectivity:4}
Dans quels cas feriez-vous tourner votre propre nœud plutôt que de recourir à un \omterm{def:nw:on-chain-connectivity:engine}{fournisseur de nœuds hébergés}~?
\end{interviewq}

\begin{interviewq}[$\star\star$ \iqroles{trader}]\label{iq:nw:on-chain-connectivity:5}
Pourquoi un trader pourrait-il envoyer une transaction en privé plutôt qu'au mempool public~?
\end{interviewq}

\begin{interviewq}[$\star\star\star$ \iqroles{développeur, chercheur}]\label{iq:nw:on-chain-connectivity:6}
Concevez la couche d'envoi des transactions d'une société de trading active sur une chaîne rapide dont les leaders se trouvent sur trois continents.
\end{interviewq}
